\chapter{Hiccups}
\index{Hiccups}

\section*{Definition and Overview}
Hiccups are sudden, involuntary contractions of the diaphragm, followed immediately by closure of the vocal cords, which produces the characteristic "hic" sound. They are usually brief and harmless, triggered by eating too quickly, drinking carbonated drinks, or excitement. Persistent hiccups lasting more than 48 hours, or intractable hiccups lasting more than a month, are uncommon and can indicate an underlying medical condition affecting the nerves that control the diaphragm, the brain, or the upper abdomen. Treatment ranges from simple home remedies to medication and treatment of the underlying cause.

\section*{Epidemiology}
Almost everyone experiences short-lived hiccups at some point, and they are more common in men than women. Persistent hiccups are rare, affecting roughly 1 in 100,000 people, and are more common in older men and in people with neurological, gastrointestinal, or metabolic conditions. Hiccups lasting days or weeks are a recognised complication of some strokes, brain tumours, and advanced cancer, and can seriously interfere with eating, sleeping, and quality of life.

\section*{Aetiology and Pathophysiology}
\begin{itemize}
    \item \textbf{Normal triggers:} eating too quickly, overeating, carbonated drinks, alcohol, sudden excitement, and temperature changes
    \item \textbf{Reflex pathway:} the hiccup reflex involves the phrenic and vagus nerves and the brainstem, which coordinate the diaphragm spasm
    \item \textbf{Gastrointestinal causes:} gastric distension, reflux, and hiatus hernia
    \item \textbf{Neurological causes:} stroke, brain tumours, head injury, multiple sclerosis, and encephalitis
    \item \textbf{Metabolic causes:} kidney failure, diabetes, electrolyte imbalance, and alcohol intoxication
    \item \textbf{Drugs and toxins:} steroids, chemotherapy, and some anaesthetic agents
    \item \textbf{Thoracic causes:} pneumonia, pleurisy, and tumours pressing on the diaphragm
    \item \textbf{Psychogenic:} stress and anxiety in some cases
\end{itemize}

\section*{Clinical Features}
\begin{itemize}
    \item Repetitive "hic" sounds at intervals of a few seconds to minutes
    \item Usually self-limiting, lasting minutes to hours
    \item Persistent hiccups last more than 48 hours
    \item Intractable hiccups last more than a month
    \item Associated symptoms reflect the underlying cause: reflux, neurological signs, or weight loss
    \item Severe or prolonged hiccups can disturb sleep, eating, and speech
    \item Fatigue and dehydration may follow prolonged episodes
\end{itemize}

\section*{Differential Diagnosis}
\begin{itemize}
    \item Myoclonus and other involuntary movements
    \item Tics and habit spasms
    \item Breathing disorders such as sighing dyspnoea
    \item The underlying cause of persistent hiccups: reflux, stroke, kidney failure, diabetes, and medications must be considered
    \item Investigation targets the underlying condition rather than the hiccup itself
\end{itemize}

\section*{When to Seek Medical Care}
See a doctor for hiccups lasting more than 48 hours, hiccups that interfere with eating or sleeping, or hiccups accompanied by other symptoms such as weight loss, fever, or neurological problems.

\textbf{Contact your local emergency services for:}
\begin{itemize}
    \item Hiccups with sudden weakness, facial drooping, or difficulty speaking (stroke)
    \item Hiccups with severe headache, confusion, or vomiting
    \item Hiccups with chest pain or severe breathlessness
    \item Collapse or loss of consciousness
\end{itemize}

\section*{Diagnosis and Investigations}
\begin{itemize}
    \item \textbf{History:} triggers, duration, medications, alcohol, and associated symptoms
    \item \textbf{Examination:} neurological, abdominal, and chest examination
    \item \textbf{Blood tests:} electrolytes, kidney function, blood sugar, and liver function
    \item \textbf{Gastroscopy:} for suspected reflux or hiatus hernia
    \item \textbf{Imaging:} chest X-ray, CT of the chest and head for structural causes
    \item \textbf{Medication review:} identify drugs that cause hiccups
    \item \textbf{Specialist referral:} to neurology, gastroenterology, or palliative care as indicated
\end{itemize}

\section*{Management}
\begin{itemize}
    \item \textbf{Home remedies:} holding the breath, drinking cold water, swallowing sugar, and other vagal manoeuvres for short episodes
    \item \textbf{Treat the cause:} stopping culprit medications, treating reflux, correcting electrolytes, and managing diabetes or kidney disease
    \item \textbf{Medication:} chlorpromazine, baclofen, gabapentin, and metoclopramide are used for persistent hiccups
    \item \textbf{Procedures:} phrenic nerve block or pacing for intractable cases
    \item \textbf{Palliative care:} important in advanced illness where hiccups add to suffering
    \item \textbf{Support:} nutritional and sleep support during prolonged episodes
\end{itemize}

\section*{Complications}
\begin{itemize}
    \item Sleep deprivation and fatigue
    \item Difficulty eating, drinking, and speaking, with weight loss and dehydration
    \item Aspiration of food or drink during episodes
    \item Postoperative hiccups after surgery or anaesthesia
    \item Depression and reduced quality of life with prolonged hiccups
    \item The underlying cause may be a serious condition requiring treatment
\end{itemize}

\section*{Prognosis and Outcomes}
Ordinary hiccups resolve spontaneously and need no treatment. Persistent hiccups usually respond to treating the underlying cause and, where needed, medication. Intractable hiccups are challenging but manageable with specialist input, and in palliative care, symptom control substantially improves comfort. The prognosis depends mainly on the underlying condition: for most people, hiccups are a passing nuisance with no lasting consequence.

\section*{Prevention and Self-Care}
\begin{itemize}
    \item Eat slowly and avoid overeating
    \item Limit carbonated drinks and alcohol
    \item Avoid very hot or very cold foods together
    \item Manage stress and excitement
    \item For short episodes, try simple measures such as sipping cold water or holding the breath
    \item Seek medical review for hiccups lasting more than 48 hours
    \item Review medications with your doctor if they trigger hiccups
    \item Treat reflux and other predisposing conditions
\end{itemize}

\section*{Sources and Further Reading}
The following reputable sources provide further information for healthcare students and patients seeking reliable, current advice about hiccups.

\begin{itemize}
    \item Healthdirect. \textit{Hiccups: causes and treatment.} https://www.healthdirect.gov.au/hiccups
    \item Better Health Channel. \textit{Hiccups.} https://www.betterhealth.vic.gov.au/
    \item NHS. \textit{Hiccups: overview.} https://www.nhs.uk/conditions/hiccups/
    \item Mayo Clinic. \textit{Hiccups: symptoms and causes.} https://www.mayoclinic.org/
    \item MSD Manual. \textit{Hiccups.} https://www.msdmanuals.com/
    \item MedlinePlus. \textit{Hiccups.} https://medlineplus.gov/ency/article/003068.htm
\end{itemize}
